\section{Flight Routes}
\label{sec:cses-flight-routes}

\problemstatement{Given a directed graph of $n$ cities and $m$ flights with
prices, find the $k$ cheapest route prices from city $1$ to city $n$.
Routes may reuse cities and flights, and different routes with equal price
are listed separately.}

\begin{patternbox}
The $k$ cheapest arrival costs is a \textbf{modified Dijkstra} in which each
node may be finalised up to $k$ times. Pop costs in increasing order; the
first $k$ times node $n$ is popped are its $k$ cheapest route prices.
\end{patternbox}

\subsection*{Let each node pop $k$ times}

Standard Dijkstra finalises each node once. Here, keep a counter of how many
times each node has been popped from the min-heap. Pop the smallest cost;
if the popped node has already been finalised $k$ times, skip it; otherwise
increment its counter and, only while the counter is $\le k$, relax its
edges (pushing $\textit{cost}+w$ for each). The $k$ costs recorded when node
$n$ is popped, in order, are the answer.

\begin{ideabox}[Bounded Multiplicity Gives k-Shortest]
Because the heap always yields the globally smallest unfinalised cost, the
$i$-th time a node is popped corresponds to its $i$-th cheapest arrival.
Capping each node at $k$ pops keeps the heap small while still producing the
$k$ best costs to the destination.
\end{ideabox}

\subsection*{Algorithm}

\begin{enumerate}
  \item Min-heap seeded with $(0,1)$; a pop-count array.
  \item Pop smallest $(d,u)$; if $u$ popped $k$ times, skip; else count it,
        record if $u=n$, and relax edges.
  \item Output the $k$ recorded costs for node $n$.
\end{enumerate}

\subsection*{Dry run}

\dryrun{$1\!\to\!2(1)$, $1\!\to\!2(3)$, $2\!\to\!3(1)$; $k=2$, $n=3$.}

\begin{itemize}
  \item Cheapest to $3$: $1{+}1=2$; second cheapest: $3{+}1=4$.
  \item Output $2\ 4$.
\end{itemize}

\subsection*{C++ implementation}

\begin{lstlisting}
#include <bits/stdc++.h>
using namespace std;

int main() {
    int n, m, k;
    cin >> n >> m >> k;
    vector<vector<pair<int,long long>>> adj(n + 1);
    for (int i = 0; i < m; i++) {
        int a, b; long long w; cin >> a >> b >> w;
        adj[a].push_back({b, w});
    }

    vector<int> cnt(n + 1, 0);
    vector<long long> best;
    priority_queue<pair<long long,int>, vector<pair<long long,int>>,
                   greater<>> pq;
    pq.push({0, 1});
    while (!pq.empty()) {
        auto [d, u] = pq.top(); pq.pop();
        if (cnt[u] >= k) continue;                 // already have k for u
        cnt[u]++;
        if (u == n) best.push_back(d);
        if (cnt[u] > k) continue;
        for (auto& [v, w] : adj[u])
            if (cnt[v] < k) pq.push({d + w, v});
    }
    for (long long c : best) cout << c << " ";
    cout << "\n";
    return 0;
}
\end{lstlisting}

\begin{complexitybox}
\textbf{Time Complexity.} $O(mk\log(mk))$ -- each edge may be pushed up to
$k$ times. \textbf{Space Complexity.} $O(mk)$ for the heap in the worst
case.
\end{complexitybox}

\begin{takeawaybox}
The $k$ shortest (repeatable) route costs come from a Dijkstra that lets
each node be finalised $k$ times. Capping pops per node bounds the work
while the heap's increasing-cost order guarantees the first $k$ arrivals at
the sink are the cheapest.
\end{takeawaybox}
